% ==========================================================================
%  Section 12 -- Statement contracts and the axiom firewall.  Owner: W10.
%  Frozen label: sec:lean-contracts.
%
%  Sources read for this section: formal/lean/challenges/_registry.yaml (the
%  two entries PPHI2_FIXED_VOLUME_COVARIANCE_LIMIT and
%  PPHI2_FIXED_VOLUME_FREE_FIELD_MEASURE_LIMIT, and permitted_axioms),
%  the per-node <NODE>.hashes.json digest files, tools/lean_challenge.py
%  (the runner), and ext/axiom_transcript.txt.
%
%  Traps observed here.  (a) CONTRACT section 11: a SKIP is not a pass, so
%  the remote_only paragraph says so in its first sentence and again at its
%  end.  (b) A footprint is a statement about dependence, never about the
%  intended statement; the last subsection keeps the two roles apart.
%  (c) The appendix labels are declared in main.tex, which the paper gate
%  does not scan for labels, so the appendices are named in words rather
%  than cross-referenced with \cref.
% ==========================================================================
\section{Statement contracts and the axiom firewall}\label{sec:lean-contracts}

Two devices carry the discipline of the formal side, and neither of them is a
proof: a \emph{statement contract}, which fixes what is to be proved, and an
\emph{axiom firewall}, which bounds what the proof of it may assume. This
section describes both, and separates what each establishes from what it leaves
open; the divergences that remain are the subject of \cref{sec:lean-boundary}.

A statement contract is an entry of the statement registry of the development,
\lmod{formal/lean/challenges/_registry.yaml}. For one node of the development
it freezes three things: the modules of the delivered package that declare the
results, the theorem names in fully qualified form, and a digest of each
statement, held beside the entry in a per-node file. The registry carries
\NumContracts{} contracts. The two this paper reports are
\lmod{PPHI2_FIXED_VOLUME_COVARIANCE_LIMIT} and
\lmod{PPHI2_FIXED_VOLUME_FREE_FIELD_MEASURE_LIMIT}. The first freezes
the covariance declarations, among them
\lean{fixedVolumeCovarianceInner_tendsto_torusFourierGreenPairing}; the second
freezes the measure-limit declarations, among them
\lean{periodicHMinusThree_fixedVolume_weak_limit_exists}. ``Fully qualified''
means that the namespace is part of the frozen string: each of the
\NumChainTheorems{} names this paper reports lives in
\lmod{Pphi2Clean.Continuum}, and the ledger of Appendix~A omits that prefix
only because the omission is itself performed mechanically.

\subsection{What freezing buys}

The digest is taken of what Lean prints for \texttt{\#check @} of a name, that
is, of the elaborated type with all implicit and instance binders, not of the
source text as typed. A later strengthening, weakening or reparametrisation of a
statement therefore produces a different digest. The runner reports the mismatch
and refuses the statement; it does not overwrite the recorded value, so the
contract stays red across runs, and the only way past it is to record the new
statement together with a written reason. That is what freezing buys: a
statement cannot change quietly, and a change accepted once cannot become
invisible afterwards.

\subsection{Both contracts are remote-only}

Both entries are marked remote-only and name the separate Lake project that hosts
the recorded digests, \lmod{/home/zhengyj/AImath-lean/AImathPPhi2Clean}. The
runner reads that field before it does anything else, and reports \lmod{SKIP} with
the declared project as the reason. \textbf{A \lmod{SKIP} is not a pass.} It
records that the digest was neither confirmed nor refuted locally, because the
recorded values live beside that project; a reader who takes the absence of a
failure for a success reads the output backwards. The flag is about \emph{where
the digest runner points}, not about whether the package can be elaborated: the
package is in this repository and, as \cref{sec:lean-overview} and the
reproduction section record, it builds in place. The evidence standing behind
these two contracts is therefore two-part: the digests and the registry entry on
the project side, and, on this side, the transcript of Appendix~B, produced by
building the modules that declare the \NumChainTheorems{} theorems and elaborating one audit file
against them.


\subsection{The axiom firewall}

The permitted set of the development is \NumPermittedAxioms{} names, the axioms
of the Lean kernel and of \lmod{Mathlib}:
\[
  \lmod{propext},\qquad \lmod{Classical.choice},\qquad \lmod{Quot.sound}.
\]
Every declaration this paper reports has
been checked with \texttt{\#print axioms}; each contract's check file issues that
command once per frozen name, next to the \texttt{\#check} that produces the
digest, so one file exercises both devices. This paper reports
\NumChainTheorems{} declarations, and the footprint of every one of them is equal
to the permitted set: no axiom introduced by the project appears, and no
placeholder standing for an omitted proof appears. Two thresholds are in force
and they are not the same one. The contract runner refuses a footprint that
contains any name outside the permitted set. The generator of Appendix~B applies
two further rules of its own: it refuses a transcript that is empty, that is not
tab-separated, whose rows are not all equal to one another, or any of whose rows
leaves the permitted set. The equality with the permitted set \emph{as a set} is
enforced one step later, by the paper's own gate, which is why a footprint that
merely stayed inside the permitted set would not reach this paper.

\subsection{The negative control}

A test that has never failed has not been shown to be able to fail. The package
therefore ships a probe, \lmod{audit/negative_axiom.lean},
\NumNegativeControlLines{} lines long: it injects
\texttt{axiom injectedAxiom : False}, derives
\texttt{theorem polluted : False := injectedAxiom}, and prints the footprint of
the derived theorem, and one line of its run output records a footprint carrying
the injected axiom and none of the permitted three. What that
establishes is that the footprint test has a case in which it must fail, so a
green run of it is not an artefact of a check that can only say yes. What it does
not establish is anything about any particular declaration: the probe is a
separate file, it is not a perturbation of the frozen names, its recorded output
is a bare measured line rather than a transcribed session with commands and a
date, and it exercises the footprint test only, not the digest. The control is
therefore attested textually -- the file is in the tree and its output line is
committed beside it -- and not by a transcript of the standard of Appendix~B.

\subsection{Two roles: dependence and intent}

A footprint is a statement about dependence, and dependence is not content.
Equality with the permitted set says that the proof term of the declaration uses
no axiom outside that set; it says nothing about whether the statement is the
intended one, because the footprint is read off the proof term and the shape of
the statement is not part of what it records. A declaration that proves a weaker
proposition, or a different one, has the same footprint as a declaration that
proves the intended proposition. What addresses intent is the other two devices.
The digest of the contract detects a change in a statement over time, and the
ledger of Appendix~A fixes, for every reported declaration, its module and its
role in the argument. The division of labour is therefore: the firewall reports
what was assumed, the digest reports what was frozen, the ledger reports where
the declaration sits. All three are necessary, none is sufficient, and what the
three together still leave unchecked is stated in \cref{sec:lean-reproduction}.
